%!TEX encoding = UTF-8 Unicode

\documentclass[t,compress,hyperref={hidelinks}]{beamer}

%include preamble.fmt

\newcommand{\lectureno}{II}

\begin{document}

%include document.fmt

\title{Type theory and logic}
\subtitle{Lecture~\lectureno: dependent type theory}
\date{2 July 2014}
\author{{\fontspec[ExternalLocation=/Library/Fonts/]{DFHeiStd-W5}\href{http://www.cs.ox.ac.uk/people/hsiang-shang.ko/}{柯向上}}}
\institute{\href{http://www.cs.ox.ac.uk/}{Department of Computer Science}\\ \href{http://www.ox.ac.uk/}{University of Oxford}\\ \href{mailto:Hsiang-Shang.Ko@@cs.ox.ac.uk}{Hsiang-Shang.Ko\char64 cs.ox.ac.uk}}

{\usebackgroundtemplate{\includegraphics[width=\paperwidth]{banner.pdf}}
\begin{frame}
\titlepage
\end{frame}}

\begin{frame}{Indexed families of sets (predicates)}

Common mathematical statements involve predicates and universal/existential quantification.

For example: ``For all natural number |x : Nat|, if |x|~is not zero, then there exists |y : Nat| such that |x|~is equal to |1 + y|.''

In type theory, a predicate on~|A| has type |A -> U| --- a \emph{family of sets} indexed by the domain~|A|.
For example:
\[ |ENTAILS \x BIND {-"\text{``if }"-} x {-"\text{\kern1pt is zero then }"-} BOT {-"\text{\kern1pt else }"-} TOP {-"\text{\kern-.1em''''}"-} : Nat -> U| \]

(Note that the above treatment is in fact unfounded in our current theory. Why? (We will fix it on Thursday.))

\end{frame}

\begin{frame}{Dependent product types (universal quantification)}

\begin{itemize}

\item Formation:
\[ \AXC{|Gamma ENTAILS A : U|} \AXC{|Gamma ENTAILS B : A -> U|}
\rule{$\Pi$F}\BIC{|Gamma ENTAILS Pi' A B : U|} \DP \]
\item Introduction:
\[ \AXC{|Gamma, x : A ENTAILS t : B x|}
\rule{$\Pi$I}\UIC{|Gamma ENTAILS \x BIND t : Pi' A B|} \DP \]
\item Elimination:
\[ \AXC{|Gamma ENTAILS f : Pi' A B|} \AXC{|Gamma ENTAILS a : A|}
\rule{$\Pi$E}\BIC{|Gamma ENTAILS f a : B a|} \DP \]

\end{itemize}

\marker{Notation} We usually write |(Pi x A) B x| for |Pi' A B| (where `|(Pi x A)|\kern-1.5pt' is regarded as a quantifier).

\marker{Exercise} Let |Gamma DEFEQ A : U, B : U, C : A -> B -> U|.
Derive
\[ |Gamma ENTAILS _ : ((Pi x A) (Pi y B) C x y) -> (Pi y B) (Pi x A) C x y| \]

\end{frame}

\begin{frame}{Dependent sum types (existential quantification)}

\begin{itemize}

\item Formation:
\[ \AXC{|Gamma ENTAILS A : U|} \AXC{|Gamma ENTAILS B : A -> U|}
\rule{$\Sigma$F}\BIC{|Gamma ENTAILS Sigma' A B : U|} \DP \]
\item Introduction:
\[ \AXC{|Gamma ENTAILS a : A|} \AXC{|Gamma ENTAILS b : B a|}
\rule{$\Sigma$I}\BIC{|Gamma ENTAILS (a, b) : Sigma' A B|} \DP \]
\item Elimination:
\[ \AXC{|Gamma ENTAILS p : Sigma' A B|}
\rule{$\Sigma$EL}\UIC{|Gamma ENTAILS fst p : A|} \DP
\quad
\AXC{|Gamma ENTAILS p : Sigma' A B|}
\rule{$\Sigma$ER}\UIC{|Gamma ENTAILS snd p : B (fst p)|} \DP \]

\end{itemize}

\marker{Notation} We usually write |(Sigma x A) B x| for |Sigma' A B|.

\marker{Exercise} Let |Gamma DEFEQ A : U, B : U, C : A -> B -> U|.
Derive the \emph{axiom of choice}:
\[ \hskip-.5em |Gamma ENTAILS _ : ((Pi x A) (Sigma y B) C x y) -> (Sigma f (A -> B)) (Pi x A) C x (f x)| \]

\end{frame}

\begin{frame}{Computation}

Let |Gamma DEFEQ A : U, B : A -> U, C : A -> U|.
Try to derive
\[ |Gamma ENTAILS _ : ((Pi p (Sigma' A B)) C (fst p)) -> ((Pi x A) B x -> C x)| \]
\ldots\ and you should notice some problems.

So far we have been concentrating on the \emph{statics} of type theory; here we need to invoke the \emph{dynamics} of the theory.

\end{frame}

\begin{frame}{Equality judgements and computation rules}

We introduce a new kind of judgements stating that two terms should be regarded as the same during typechecking:
\[ |Gamma ENTAILS t = u IN A| \]

For each set, (when applicable) we specify additional \emph{computational rules} stating that eliminating an introductory term yields a component of the latter.
For example, for product types we have two computation rules:
\vskip-3ex
\[ \kern-.75em
\AXC{|Gamma ENTAILS a : A|} \AXC{|Gamma ENTAILS b : B|}
\rule{$\times$CL}\insertBetweenHyps{\hskip.75em}\BIC{|Gamma ENTAILS fst (a , b) = a IN A|} \DP
\quad
\AXC{|Gamma ENTAILS a : A|} \AXC{|Gamma ENTAILS b : B|}
\rule{$\times$CR}\insertBetweenHyps{\hskip.75em}\BIC{|Gamma ENTAILS snd (a , b) = b IN B|} \DP \]

\end{frame}

\begin{frame}{More computation rules}

\[ \AXC{|Gamma, x : A ENTAILS t : B|} \AXC{|Gamma ENTAILS a : A|}
\rule{$\to$C}\BIC{|Gamma ENTAILS (\x BIND t) a = t(subst a x) IN A -> B|} \DP \]


\[ \AXC{|Gamma ENTAILS a : A|} \AXC{|Gamma ENTAILS f : A -> C|} \AXC{|Gamma ENTAILS g : B -> C|}
\rule{$+$CL}\TIC{|Gamma ENTAILS case (left a) f g = f a IN C|} \DP \]

\[ \AXC{|Gamma ENTAILS b : B|} \AXC{|Gamma ENTAILS f : A -> C|} \AXC{|Gamma ENTAILS g : B -> C|}
\rule{$+$CR}\TIC{|Gamma ENTAILS case (right b) f g = g b IN C|} \DP \]

\end{frame}


\begin{frame}{More computation rules}

\[ \AXC{|Gamma, x : A ENTAILS t : B x|} \AXC{|Gamma ENTAILS a : A|}
\rule{$\Pi$C}\BIC{|Gamma ENTAILS (\x BIND t) a = t(subst a x) IN B a|} \DP \]

\[ \AXC{|Gamma ENTAILS a : A|} \AXC{|Gamma ENTAILS b : B a|}
\rule{$\Sigma$CL}\BIC{|Gamma ENTAILS fst (a , b) = a IN A|} \DP \]

\[ \AXC{|Gamma ENTAILS a : A|} \AXC{|Gamma ENTAILS b : B a|}
\rule{$\Sigma$CR}\BIC{|Gamma ENTAILS snd (a , b) = b IN B a|} \DP \]

\end{frame}

\begin{frame}
Predicates respect computation.
\end{frame}

\end{document}
